% ====================================================================== Task 2A
\clearpage
\section{Task 2A: Instruction-Conscious L2 Replacement (I-RRIP)}

\subsection{Why protect instruction lines in the L2?}
When a load misses in the L1D and hits in the L2, it waits about 10 extra cycles, and the out-of-order core
(352-entry ROB) can usually hide much of that by working on other instructions. An L1I miss is different: fetch
stops until the line arrives, so the delay goes almost straight into execution time.

Our traces have large code footprints. Even with SRRIP the L2 instruction load MPKI is 10.55, 19.32 and 1.20, and
on Trace2 instruction misses are 71\% of all L2 load misses. Keeping instruction lines in the L2 turns these
expensive front-end stalls into short L2 hits. The price is some extra data misses, which the core partly hides anyway.

\subsection{Design}
I-RRIP keeps the SRRIP hardware: a 2-bit re-reference prediction value (RRPV) per line, the victim is a line with
RRPV\,3, and if no line has RRPV\,3 every line of the set ages by one. I-RRIP changes all three parts of the policy:

\begin{table}[H]\centering\small
\begin{tabularx}{\textwidth}{@{}l>{\raggedright\arraybackslash}p{40mm}>{\raggedright\arraybackslash}X@{}}
\toprule
Part & SRRIP & I-RRIP \\
\midrule
Insertion & every line at RRPV 2               & instruction line at RRPV \textbf{0}, data line at RRPV \textbf{3} \\
Promotion & every hit $\to$ RRPV 0             & demand hits (instruction and data) $\to$ RRPV 0; \textbf{a writeback hit does not promote} \\
Eviction  & first line (lowest way) with RRPV 3 & among lines with RRPV 3, \textbf{a data line before an instruction line} \\
\bottomrule
\end{tabularx}
\end{table}

\textbf{Hardware cost over SRRIP:} one instruction bit per line, i.e.\ 8192 bits for the 512\,KB L2 (ChampSim already
keeps it in \code{BLOCK::instruction}). Everything else is just different constants and one extra comparison in
the victim search.

\subsection{How each part was chosen: ablation}
Starting from SRRIP, we changed one part at a time and kept a change only if it raised the geometric-mean speedup
by more than 0.05\%. Smaller gaps count as ties, the same rule as in Task 1. Table~\ref{tab:t2a-abl} lists every
variant and Figure~\ref{fig:t2a-abl} plots their geometric-mean speedups.

\begin{table}[H]\centering\small
\setlength{\tabcolsep}{3.6pt}
\begin{tabular}{@{}lccccccccccc@{}}
\toprule
Name & \makecell{Instr.\\insert} & \makecell{Data\\insert} & \makecell{Data\\hit} & \makecell{WB\\hit} & Victim
     & \makecell{Speedup\\T1} & \makecell{Speedup\\T2} & \makecell{Speedup\\T3} & Geomean
     & \makecell{L2 instr.\\MPKI (mean)} & \makecell{L2 data\\MPKI (mean)} \\
\midrule
SRRIP & 2 & 2 & $\to$0 & $\to$0 & first at 3 & 0.9996 & 1.0020 & 0.9990 & 1.0002 & 10.36 & 5.95 \\
A1    & 1 & 2 & $\to$0 & $\to$0 & first at 3 & 1.0004 & 1.0064 & 0.9996 & 1.0021 & 9.90 & 6.18 \\
A2    & 0 & 2 & $\to$0 & $\to$0 & first at 3 & 0.9996 & 1.0097 & 0.9998 & 1.0030 & 9.65 & 6.37 \\
A3    & 1 & 3 & $\to$0 & $\to$0 & first at 3 & 0.9974 & 1.0201 & 0.9998 & 1.0057 & 8.19 & 7.26 \\
A4    & 0 & 3 & $\to$0 & $\to$0 & first at 3 & 0.9970 & 1.0216 & 1.0000 & 1.0061 & 8.07 & 7.39 \\
P1    & 0 & 2 & $-1$   & $\to$0 & first at 3 & 0.9992 & 1.0136 & 1.0001 & 1.0043 & 9.18 & 6.55 \\
P2    & 0 & 3 & $-1$   & $\to$0 & first at 3 & 0.9964 & 1.0238 & 1.0015 & 1.0072 & 7.70 & 7.61 \\
P3    & 0 & 2 & $\to$0 & none   & first at 3 & 0.9990 & 1.0138 & 1.0004 & 1.0044 & 9.16 & 6.44 \\
P4    & 0 & 3 & $\to$0 & none   & first at 3 & 0.9974 & 1.0267 & 1.0046 & 1.0095 & 7.71 & 7.56 \\
P5    & 0 & 3 & $-1$   & none   & first at 3 & 0.9963 & 1.0269 & 1.0059 & 1.0096 & 7.50 & 7.80 \\
E1    & 0 & 2 & $\to$0 & $\to$0 & data first & 0.9989 & 1.0162 & 0.9999 & 1.0050 & 9.01 & 6.57 \\
E2    & 0 & 3 & $\to$0 & $\to$0 & data first & 0.9961 & 1.0225 & 1.0000 & 1.0061 & 7.70 & 7.71 \\
E3    & 0 & 2 & $\to$0 & none   & data first & 0.9987 & 1.0188 & 1.0004 & 1.0059 & 8.65 & 6.62 \\
\chosen{I-RRIP} & 0 & 3 & $\to$0 & none & data first & 0.9962 & 1.0276 & 1.0050 & 1.0095 & 7.44 & 7.91 \\
E4    & 0 & 3 & $-1$   & none   & data first & 0.9956 & 1.0271 & 1.0061 & 1.0095 & 7.44 & 7.91 \\
\bottomrule
\end{tabular}
\caption{I-RRIP ablation. The insert columns are RRPVs. ``Data hit'' $\to$0 means promotion to RRPV 0, $-1$ means one
step closer. ``WB hit'' is the promotion on a writeback hit. Speedup is over LRU. Highlighted: the submitted I-RRIP.}
\label{tab:t2a-abl}
\end{table}

\begin{figure}[H]
\centering
\begin{tikzpicture}
\begin{axis}[base, xbar, width=0.8\textwidth, height=72mm, bar width=6.5pt,
  ymajorgrids=false, xmajorgrids,
  y dir=reverse, ytick={0,...,14},
  yticklabels={SRRIP,A1,A2,A3,A4,P1,P2,P3,P4,P5,E1,E2,E3,I-RRIP,E4},
  ytick style={draw=none}, enlarge y limits={abs=0.6},
  xmin=0.998, xmax=1.0112, xtick distance=0.002,
  xticklabel style={/pgf/number format/.cd, fixed, fixed zerofill, precision=3},
  xlabel={Geometric-mean speedup over LRU (3 traces)},
  nodes near coords, nodes near coords align={horizontal},
  every node near coord/.append style={font=\scriptsize, /pgf/number format/.cd, fixed, fixed zerofill, precision=4},
  clip=false]
  \addplot[bBlue, bar shift=0pt] coordinates {
    (1.0002,0) (1.0021,1) (1.0030,2) (1.0057,3) (1.0061,4) (1.0043,5) (1.0072,6)
    (1.0044,7) (1.0095,8) (1.0096,9) (1.0050,10) (1.0061,11) (1.0059,12) (1.0095,14)};
  \addplot[bOrange, bar shift=0pt] coordinates {(1.0095,13)};
  \draw[black!75, dashed, thin] (axis cs:1,14.6) -- (axis cs:1,-0.6) node[above, font=\scriptsize, black!75] {LRU};
\end{axis}
\end{tikzpicture}
\caption{Geometric-mean speedup over LRU of every ablation variant (dashed line: LRU). Orange: the submitted I-RRIP.}
\label{fig:t2a-abl}
\end{figure}

\textbf{Insertion.} Inserting instruction lines at RRPV 0 instead of 2 helps (A2: $+0.30\%$), because they survive
two more aging rounds before they can be evicted. RRPV 1 protects them less and is worse (A1: $+0.21\%$). The bigger
step is inserting data at RRPV 3 (A4: $+0.61\%$, A3: $+0.57\%$): a new data line then becomes the first victim of its
set unless it is reused before the next miss. This is a safe bet, because most L2 data lines are never reused
before eviction (the L2 load hit rate of SRRIP is only 30\%, 49\% and 27\%). Little useful data is lost and
the space goes to instructions. We keep both changes: instructions at 0, data at 3.

\textbf{Promotion.} We tested two changes. Promoting a data hit by only one step (P1, P2) helps a little on its own
(A4 $\to$ P2: $+0.61\% \to +0.72\%$). The most important single change, however, is \emph{not promoting on a
writeback hit} (A4 $\to$ P4: $+0.61\% \to +0.95\%$). Writebacks are dirty lines that the L1D has just evicted. They
are 11--25\% of all L2 accesses, and over 92\% of them hit in the L2 because the line is still there from its fill.
SRRIP moves each such line to RRPV 0, so the line the L1D just threw away gets the highest priority in the L2, ahead
of every instruction line. A writeback is not a reuse by the core, so it should not change the prediction.
Once writebacks stop promoting, one-step data promotion adds nothing (P4 $+0.95\%$ vs P5 $+0.96\%$, a tie), so
demand hits keep the simpler SRRIP rule ($\to$ RRPV 0): a data line that really is reused deserves to stay.

\textbf{Eviction.} Picking a data line before an instruction line among the RRPV-3 candidates matters when both kinds
reach RRPV 3 together. This is common when data is inserted at 2 (A2 $\to$ E1: $+0.30\% \to +0.50\%$). With data
inserted at 3, most RRPV-3 lines are data lines already, and the speedup gain is a tie (P4 $+0.95\%$ vs I-RRIP
$+0.95\%$). We keep the rule anyway: it is the only part that protects an \emph{old} instruction line, and it
lowers the L2 instruction MPKI further, which is the stated goal of the task (Trace2: 12.98 $\to$ 12.33). It costs
nothing beyond the instruction bit.

\textbf{Final choice.} I-RRIP ties for the highest geometric-mean speedup with P4, P5 and E4 (all within 0.01\%), and
together with E4 it has the lowest mean L2 instruction MPKI. E4 only adds one-step data promotion, which brings no
gain, so we submit the simpler I-RRIP. The cost is visible on Trace1, the trace whose data is reused in the L2:
its data MPKI rises from 7.06 to 9.63 and its speedup drops to 0.9962. Task 2B fixes this.

\subsection{Results}
\begin{table}[H]\centering\small
\begin{tabular}{@{}lccccc@{}}
\toprule
Configuration & IPC & Speedup & L2C load MPKI & L2 instr.\ load MPKI & L2 data load MPKI \\
\midrule
\grp{6}{Trace1}
LRU (baseline)   & 0.2321 & 1.0000 & 17.70 & 10.37 & 7.33 \\
SRRIP            & 0.2320 & 0.9996 & 17.61 & 10.55 & 7.06 \\
\chosen{I-RRIP}  & 0.2312 & 0.9962 & 18.88 & 9.25  & 9.63 \\
\midrule
\grp{6}{Trace2}
LRU (baseline)   & 0.2499 & 1.0000 & 28.64 & 20.14 & 8.50 \\
SRRIP            & 0.2504 & 1.0020 & 27.27 & 19.32 & 7.95 \\
\chosen{I-RRIP}  & 0.2568 & 1.0276 & 23.25 & 12.33 & 10.92 \\
\midrule
\grp{6}{Trace3}
LRU (baseline)   & 0.3820 & 1.0000 & 3.96 & 1.12 & 2.84 \\
SRRIP            & 0.3816 & 0.9990 & 4.04 & 1.20 & 2.85 \\
\chosen{I-RRIP}  & 0.3839 & 1.0050 & 3.90 & 0.73 & 3.17 \\
\midrule
\grp{6}{Average (geometric mean for IPC and speedup, arithmetic mean otherwise)}
LRU (baseline)   & 0.2809 & 1.0000 & 16.77 & 10.54 & 6.22 \\
SRRIP            & 0.2809 & 1.0002 & 16.31 & 10.36 & 5.95 \\
\chosen{I-RRIP}  & 0.2836 & 1.0095 & 15.34 & 7.44  & 7.91 \\
\bottomrule
\end{tabular}
\caption{Task 2A: IPC, speedup over LRU and L2 load MPKIs (no prefetcher).}
\label{tab:t2a-res}
\end{table}

\replfig{lru,srrip,irrip}{LRU (baseline),SRRIP,I-RRIP}{leg:t2a}{5pt}
  {Task 2A: speedup and L2C / L2 instruction / L2 data load MPKI of LRU, SRRIP and I-RRIP.}{fig:t2a-res}

I-RRIP lowers the mean L2 instruction load MPKI from 10.36 (SRRIP) to 7.44, i.e.\ 28\% fewer instruction misses,
while the data load MPKI rises from 5.95 to 7.91. On Trace2, where instructions cause most L2 misses, this gives
a $+2.76\%$ speedup over LRU, and the total L2 load MPKI falls as well. On Trace3 the instruction misses almost
disappear (1.20 $\to$ 0.73). On Trace1 the extra data misses outweigh the saved instruction misses.
